\documentclass[../DoAn.tex]{subfiles}
\begin{document}

Chương này trình bày tổng quan về đề tài, bao gồm bối cảnh và vấn đề thực tiễn dẫn đến sự cần thiết của hệ thống, mục tiêu và phạm vi của đồ án, định hướng giải pháp được lựa chọn, cũng như bố cục của toàn bộ báo cáo.

\section{Đặt vấn đề}
\label{section:1.1}
Trong bối cảnh hội nhập kinh tế quốc tế, các doanh nghiệp ngày càng thường xuyên làm việc với đối tác, khách hàng và chi nhánh ở nhiều quốc gia khác nhau. Điều này làm cho nhu cầu xử lý, trao đổi và dịch tài liệu bằng nhiều ngôn ngữ trở nên phổ biến hơn. Các tài liệu trong doanh nghiệp không chỉ là văn bản thuần mà còn tồn tại dưới nhiều định dạng khác nhau như DOCX, XLSX, PPTX, PDF hoặc hình ảnh.

Vấn đề thứ nhất là sự thiếu nhất quán trong thuật ngữ chuyên ngành. khi sử dụng các công cụ dịch rời rạc hoặc dịch thủ công, doanh nghiệp khó kiểm soát việc cùng một thuật ngữ được dịch nhất quán giữa các tài liệu, phòng ban hay dự án khác nhau, ảnh hưởng trực tiếp đến tính chính xác và chuyên nghiệp của nội dung. Vấn đề thứ hai là việc mất cấu trúc và định dạng tài liệu sau khi dịch: với các tài liệu có bảng biểu, hình ảnh, bố cục phức tạp hoặc nhiều thành phần trình bày, phần lớn công cụ dịch hiện tại chỉ trả về văn bản thuần, buộc người dùng phải tốn thêm thời gian chỉnh sửa lại định dạng thủ công sau khi dịch.

Từ những vấn đề trên, bài toán đặt ra là cần xây dựng một hệ thống hỗ trợ dịch tài liệu đa ngôn ngữ tích hợp, vừa đảm bảo tính nhất quán thuật ngữ, vừa giữ nguyên cấu trúc và định dạng tài liệu gốc, vừa tập trung toàn bộ quy trình trong một nền tảng duy nhất. Nếu giải quyết được bài toán này, doanh nghiệp có thể nâng cao hiệu quả xử lý tài liệu, giảm thao tác thủ công và cải thiện chất lượng trao đổi thông tin đa ngôn ngữ. Ngoài lĩnh vực doanh nghiệp, bài toán này cũng có thể được áp dụng trong giáo dục, thương mại, pháp lý và nhiều lĩnh vực khác có nhu cầu xử lý tài liệu đa ngôn ngữ.

\section{Mục tiêu và phạm vi đề tài}
\label{section:1.2}
Hiện nay, nhu cầu dịch tài liệu đa ngôn ngữ trong doanh nghiệp ngày càng trở nên phổ biến, đặc biệt với các tài liệu phục vụ trao đổi nội bộ, làm việc với đối tác, khách hàng hoặc chi nhánh ở nhiều quốc gia. Một số cách tiếp cận thường được sử dụng là dịch thủ công, sử dụng công cụ dịch trực tuyến hoặc dùng các phần mềm hỗ trợ dịch tài liệu. Các công cụ này có ưu điểm là dễ tiếp cận, hỗ trợ nhiều ngôn ngữ và có thể xử lý nhanh các đoạn văn bản ngắn.

Tuy nhiên, khi áp dụng vào tài liệu văn phòng trong môi trường doanh nghiệp, các cách tiếp cận trên vẫn còn nhiều hạn chế. Phần lớn công cụ dịch thông thường chỉ phù hợp với văn bản thuần, trong khi tài liệu thực tế thường tồn tại dưới nhiều định dạng như DOCX, XLSX, PPTX, PDF hoặc hình ảnh. Với các tài liệu có bảng biểu, bố cục phức tạp, slide trình bày hoặc nội dung được trích xuất từ ảnh, việc dịch xong nhưng vẫn giữ nguyên cấu trúc và định dạng ban đầu là một thách thức lớn. Người dùng thường phải chỉnh sửa lại thủ công sau khi dịch, gây mất thời gian và dễ phát sinh sai sót.

Bên cạnh đó, một vấn đề quan trọng khác là tính thống nhất của thuật ngữ chuyên ngành. Trong doanh nghiệp, cùng một thuật ngữ có thể xuất hiện trong nhiều tài liệu khác nhau và cần được dịch nhất quán. Nếu chỉ sử dụng công cụ dịch rời rạc, kết quả dịch có thể thay đổi giữa các lần dịch hoặc giữa các người dùng khác nhau, làm giảm tính chuyên nghiệp và gây khó khăn trong việc sử dụng tài liệu. Vì vậy, bài toán dịch tài liệu không chỉ dừng lại ở việc chuyển đổi ngôn ngữ, mà còn cần đảm bảo khả năng kiểm soát thuật ngữ và chất lượng nội dung dịch.

Từ những hạn chế trên, đồ án hướng tới xây dựng một hệ thống hỗ trợ dịch đa ngôn ngữ cho văn bản và tài liệu văn phòng, tập trung vào việc xử lý nhiều định dạng tệp, giữ nguyên cấu trúc tài liệu sau khi dịch và hỗ trợ sử dụng thư viện thuật ngữ trong quá trình dịch. Hệ thống cung cấp các chức năng chính như dịch văn bản, dịch tệp tài liệu, xử lý các định dạng DOCX, XLSX, PPTX, PDF, ảnh và OCR, cho phép lựa chọn chế độ sử dụng thư viện thuật ngữ chung, thư viện thuật ngữ riêng hoặc không dùng thư viện thuật ngữ, đồng thời lưu lại kết quả dịch để người dùng có thể tải về và quản lý.

Mục tiêu của đồ án là tích hợp quá trình dịch nội dung, xử lý định dạng tài liệu và quản lý thuật ngữ trong cùng một hệ thống. Nhờ đó, người dùng có thể dịch tài liệu đa ngôn ngữ thuận tiện hơn, giảm bớt thao tác chỉnh sửa thủ công sau khi dịch và cải thiện tính nhất quán của thuật ngữ chuyên ngành. Đây là hướng tiếp cận phù hợp với nhu cầu thực tế của doanh nghiệp, đặc biệt trong các môi trường thường xuyên xử lý khối lượng lớn tài liệu bằng nhiều ngôn ngữ khác nhau.


\section{Định hướng giải pháp}
\label{section:1.3}
Để giải quyết bài toán đã đặt ra, đồ án lựa chọn định hướng phát triển một hệ thống web hỗ trợ dịch đa ngôn ngữ cho văn bản và tài liệu văn phòng trong môi trường doanh nghiệp. Hệ thống được xây dựng theo mô hình Client-Server, trong đó React/Vite được sử dụng cho phần giao diện người dùng nhằm đảm bảo khả năng tương tác thuận tiện, Django REST Framework được sử dụng cho backend để xử lý nghiệp vụ, quản lý dữ liệu và cung cấp API. Đối với chức năng dịch, hệ thống tích hợp Google Cloud Translation, Google Cloud Storage cho thư viện thuật ngữ, Amazon S3 cho lưu trữ tệp và ConvertAPI cho chuyển đổi định dạng tài liệu.

Hệ thống cho phép người dùng dịch văn bản hoặc tải lên các tệp tài liệu để dịch sang nhiều ngôn ngữ khác nhau. Các định dạng được hỗ trợ bao gồm DOCX, XLSX, PPTX, PDF, ảnh và nội dung cần xử lý OCR. Bên cạnh việc dịch nội dung, hệ thống còn hướng tới việc giữ lại cấu trúc và định dạng của tài liệu gốc sau khi dịch, đồng thời cho phép người dùng lựa chọn sử dụng thư viện thuật ngữ chung, thư viện thuật ngữ riêng hoặc không dùng thư viện thuật ngữ tùy theo nhu cầu. Ngoài các chức năng dịch, hệ thống còn hỗ trợ quản lý lịch sử tệp đã dịch, quản lý người dùng và quản lý thư viện thuật ngữ. 

\section{Đóng góp của đồ án}
\label{section:1.4}

Đóng góp chính của đồ án là xây dựng được một hệ thống web tích hợp quá trình dịch tài liệu đa định dạng, xử lý định dạng và quản lý thuật ngữ trong cùng một nền tảng. Cụ thể, đồ án đóng góp các điểm sau:

\begin{itemize}
    \item Xây dựng quy trình dịch thống nhất hỗ trợ nhiều định dạng tài liệu văn phòng phổ biến gồm DOCX, XLSX, PPTX, PDF, ảnh và OCR trong một API duy nhất.
    \item Giữ nguyên cấu trúc và định dạng tài liệu gốc sau khi dịch, giảm thiểu thao tác chỉnh sửa thủ công của người dùng.
    \item Xây dựng cơ chế thư viện thuật ngữ hai cấp (chung và riêng) tích hợp trực tiếp vào quá trình dịch, cải thiện tính nhất quán của thuật ngữ chuyên ngành.
    \item Xây dựng hệ thống quản lý người dùng, phân quyền theo vai trò và quản lý lịch sử tệp đã dịch phục vụ môi trường doanh nghiệp.
\end{itemize}

Kết quả đạt được là hệ thống hoàn thiện các chức năng cốt lõi trong phạm vi thử nghiệm, có khả năng hỗ trợ doanh nghiệp dịch tài liệu nhanh hơn và quản lý thuật ngữ nhất quán hơn so với việc sử dụng các công cụ dịch rời rạc.

\section{Bố cục đồ án}
\label{section:1.5}

Phần còn lại của báo cáo đồ án tốt nghiệp được tổ chức như sau.

Chương~\ref{chapter:Related_works} tập trung khảo sát thực tế nhu cầu dịch văn bản và tài liệu văn phòng đa ngôn ngữ trong môi trường doanh nghiệp, đồng thời phân tích một số công cụ và hệ thống dịch hiện có để rút ra điểm mạnh, điểm hạn chế. Qua việc xem xét nhu cầu sử dụng của người dùng và tham khảo các giải pháp dịch tài liệu hiện nay, chương xác định các yêu cầu cần thiết đối với hệ thống cần xây dựng, đặc biệt là khả năng dịch nhiều định dạng tài liệu, giữ nguyên cấu trúc sau khi dịch và hỗ trợ thống nhất thuật ngữ chuyên ngành. Ngoài ra, chương cũng trình bày các nhóm người dùng, các use case chính, đặc tả yêu cầu chức năng và yêu cầu phi chức năng, làm cơ sở cho quá trình thiết kế và triển khai hệ thống trong các chương tiếp theo.

Chương~\ref{chapter:Methodology} trình bày cơ sở lý thuyết và các công nghệ được lựa chọn để xây dựng hệ thống. Nội dung chương bao gồm mô hình ứng dụng web Client-Server, kiến trúc REST API, cơ chế xác thực và phân quyền người dùng, cũng như các công nghệ xây dựng giao diện và backend như React/Vite và Django REST Framework. Chương cũng giới thiệu các dịch vụ và công cụ phục vụ trực tiếp cho bài toán dịch tài liệu, gồm Google Cloud Translation cho dịch máy, thư viện thuật ngữ trên Google Cloud Storage để hỗ trợ thống nhất thuật ngữ, Amazon S3 để lưu trữ tệp, ConvertAPI để chuyển đổi định dạng PDF và các thư viện xử lý tài liệu DOCX, XLSX, PPTX, PDF, ảnh và OCR. Ngoài ra, chương trình bày Docker và Nginx như các công cụ hỗ trợ đóng gói, triển khai và vận hành hệ thống.

Chương~\ref{chapter:Experiment} trình bày quá trình phân tích, thiết kế và xây dựng hệ thống dịch đa ngôn ngữ. Nội dung chương bao gồm thiết kế kiến trúc tổng thể, thiết kế cơ sở dữ liệu, xây dựng các chức năng chính như dịch văn bản, dịch tệp, quản lý lịch sử dịch, quản lý thuật ngữ và quản lý người dùng. Bên cạnh đó, chương cũng trình bày quá trình tích hợp các dịch vụ bên ngoài phục vụ dịch máy, lưu trữ và chuyển đổi tài liệu, đồng thời đánh giá kết quả kiểm thử các chức năng cốt lõi sau khi triển khai.

Chương~\ref{chapter:SolutionAndContribution} trình bày các đóng góp chính của đồ án trong việc xây dựng hệ thống dịch tài liệu đa ngôn ngữ. Nội dung chương tập trung vào các điểm nổi bật như hỗ trợ dịch tài liệu ở nhiều định dạng khác nhau, giữ lại cấu trúc tài liệu sau khi dịch và sử dụng thư viện thuật ngữ để nâng cao tính nhất quán của nội dung dịch. Qua đó, chương làm rõ các giải pháp được triển khai và giá trị của chúng đối với bài toán dịch tài liệu đa ngôn ngữ.

Chương~\ref{chapter:conclusion} tổng kết quá trình thực hiện đồ án tốt nghiệp. Chương này đánh giá những kết quả đã đạt được. Đồng thời, chương cũng chỉ ra một số hạn chế còn tồn tại như phạm vi kiểm thử chưa bao phủ hết dữ liệu người dùng thực tế, hiệu năng với khối lượng tệp lớn cần được đánh giá thêm và chất lượng dịch vẫn phụ thuộc vào dịch vụ dịch máy bên ngoài. Cuối cùng, chương đề xuất các hướng phát triển trong tương lai như mở rộng định dạng tài liệu hỗ trợ, cải thiện khả năng giữ định dạng, tối ưu tốc độ xử lý tệp, bổ sung thống kê sử dụng hệ thống và nâng cao quy trình quản trị thuật ngữ.

\end{document}
